\chapter{Capstone: Extending and Shipping a C Game}
\label{ch:ccap}

This chapter closes the C track the way Chapter~\ref{ch:asmcap} closed the
assembly track: by extending a game and shipping it. You will embed a C game in
the firmware for quick iteration, build it into a \fname{.prg32} cartridge, and
deploy it to both QEMU and the board. Along the way you will meet two ideas that
recur in serious systems programming---reusing helper functions, and doing
arithmetic without a floating-point unit---both grounded in the platform's richer
C examples \citep{prg32tutorialc,prg32repo}.

\section{Embedding a C game in the firmware}

For fast iteration during development, the platform lets you compile a C game
directly into the firmware instead of packaging a cartridge each time. You add the
game's source to the firmware's build list and call its three functions from the
main loop \citep{prg32tutorialc}. This is a temporary lab arrangement, restored
afterwards, but it shortens the edit-build-run cycle while you are still changing
the game often.

\begin{lstlisting}[style=c,caption={Calling a C game's three functions from the firmware's main loop. \texttt{init} once, then \texttt{update} and \texttt{draw} each frame, with a present and a short delay.}]
#include "prg32.h"
#include "freertos/FreeRTOS.h"
#include "freertos/task.h"

void pong_c_init(void);
void pong_c_update(void);
void pong_c_draw(void);

void app_main(void) {
    prg32_init();
    pong_c_init();
    while (1) {
        pong_c_update();
        pong_c_draw();
        prg32_gfx_present();             /* show the finished frame */
        vTaskDelay(pdMS_TO_TICKS(33));   /* ~30 frames per second */
    }
}
\end{lstlisting}

\begin{prgconcept}
The frame loop you met as a diagram in Chapter~\ref{ch:platform} is, in the
firmware, a literal \code{while (1)} that calls \code{update}, then \code{draw},
then \code{prg32\_gfx\_present} to display the frame, then waits briefly so the
game runs at a steady rate. Embedding your game here lets you see the loop you have
been writing the body of all along.
\end{prgconcept}

\begin{prgwarn}
The embedding is temporary. The platform's tutorial is explicit that you should
restore the firmware's default main loop and build list after the lab, so the
resident firmware returns to loading cartridges normally \citep{prg32tutorialc}.
Keep a clean copy, or use version control (Appendix~\ref{app:github}) so restoring
is one command.
\end{prgwarn}

\section{Reusing helpers and reading richer examples}

Once your game grows, repeated code is a liability. C's answer is the helper
function: a named, reusable piece of logic you write once and call many times. You
already saw the platformer's \code{put\_run} and \code{put\_column} helpers in
Chapter~\ref{ch:c2}; they exist precisely so the world-building code is not a wall
of near-identical lines. When you extend your own game, watch for three or more
lines that recur with small variations, and lift them into a helper.

The platform's raycaster example is the next step for a class ready for more
\citep{prg32teaching}. It renders a Doom-style corridor---a convincing pseudo-3-D
view---on the same \riscv runtime used for everything else. What makes it
instructive is not the 3-D effect but \emph{how} it computes that effect without a
floating-point unit.

\section{Arithmetic without floating point: fixed point}

Small microcontrollers often have no hardware for fractional numbers, yet a
raycaster needs fractions: half a step, a third of an angle. The classic solution,
used by the platform's raycaster, is \emph{fixed-point arithmetic}: represent a
fractional value as an integer that has been scaled up by a fixed factor, do the
arithmetic in integers, then scale back down \citep{prg32tutorialc,prg32repo}.

\begin{lstlisting}[style=c,caption={Fixed-point in the raycaster. Direction vectors are stored scaled by 256; multiplying and then dividing by 256 keeps the fraction.}]
/* dir_q8 holds direction components scaled by 256 (called Q8 fixed point). */
int hx = player_x + (dir_q8[angle][0] * dist) / 256;
int hy = player_y + (dir_q8[angle][1] * dist) / 256;
\end{lstlisting}

\begin{prgconcept}
\textbf{Fixed-point arithmetic} stores a fractional number as an integer scaled by
a fixed factor---here, 256. A value of $1.5$ is stored as $384$, because
$1.5\times256=384$. Multiplications are done in integers and then divided by the
scale to bring the result back. This lets a processor with no floating-point unit
compute smooth motion and angles using only integer instructions---a technique as
old as real-time graphics and still used in embedded systems today.
\end{prgconcept}

\begin{prgaside}
The notation ``Q8'' in the example names the format: eight bits of fraction, so a
scale of $2^8 = 256$. Fixed-point was the only practical way to do fast fractional
maths on the home computers of the retro era, and the raycaster's corridor is a
direct descendant of the techniques that made early 3-D games run on machines with
no floating-point hardware at all. Learning it here is learning a genuine piece of
computing history that remains useful.
\end{prgaside}

\section{Building and shipping a C cartridge}

Packaging a C game is identical to packaging an assembly game, except the source
is a \code{.c} file and the entry prefix matches the C function names. The same
\fname{prg32\_game.py} tool resolves the framework symbols from the firmware ELF
and produces a \fname{.prg32} cartridge \citep{prg32tutorialc}.

\begin{lstlisting}[style=bash,caption={Building a C cartridge against the board firmware.}]
python3 tools/prg32_game.py build \
  examples/games/pong/c/game.c \
  --firmware-elf build-esp32c6/PRG32.elf \
  --entry-prefix pong_c \
  --name pong-c \
  --out build-esp32c6/pong-c.prg32
\end{lstlisting}

Deployment is the same as in Chapter~\ref{ch:asmcap}: upload over the board's Wi-Fi
for hardware, or build against the QEMU firmware ELF and stage into the flash image
for the emulator \citep{prg32qemu}.

\begin{lstlisting}[style=bash,caption={Uploading the C cartridge to the board, then the QEMU variant.}]
python3 tools/prg32_game.py upload build-esp32c6/pong-c.prg32 --url http://192.168.4.1
# For QEMU: build against build-qemu/PRG32.elf, then:
python3 tools/prg32_game.py upload-qemu build-qemu/pong-c.prg32
\end{lstlisting}

\begin{prgtargets}
\textbf{Both targets.} A C cartridge and an assembly cartridge are the same kind of
file and run through the same loader; the framework neither knows nor cares which
language produced the three functions inside \citep{prg32teaching}. This is the
deepest demonstration of the platform's design: two courses, two languages, one
runtime, two targets, one cartridge format.
\end{prgtargets}

\begin{prgtry}
Take the C Pong you wrote in Chapter~\ref{ch:c1}, add a score that increments when
the ball bounces off the paddle---guarding against counting one hit many times, as
warned in Chapter~\ref{ch:asmcap}---and display it with \code{prg32\_gfx\_text8}.
Build it as a cartridge for QEMU and run it. Then, if you have a board, upload the
same cartridge over Wi-Fi and play it. You will have written, extended, packaged,
and shipped a complete C game on real \riscv hardware \citep{prg32repo}.
\end{prgtry}

\begin{prgcheck}
The C track is complete. You can embed a game for fast iteration, factor repeated
logic into helpers, reason about fixed-point arithmetic, and build and deploy a C
cartridge to both targets. With both tracks behind you (or your own track behind
you), Chapter~\ref{ch:join} offers a short, illuminating reunion: the same game,
in both languages, side by side.
\end{prgcheck}
